% Section: §1 Theory-Laden Perception as Central Phenomenological Target
% PDF context: Kuhn's three faces of paradigm-shaped cognition (what is seen, what is investigated, what alternatives are entertained) treated as expressions of a single structural feature, defended as a stronger rhetorical commitment than an accumulated-checklist account.
% Drawing from: cluster_C6_theory_laden
% Entries cited: 9
% Full-text papers consulted: 0 (only available GROBID text was a book review of Sarkar2006_11, not substantive)
% Key renaming TODOs:
%   Chang2012_9   -> chang2012WaterChemRev      % Anchor: Chang on phlogiston/water, foundational chemistry-revolution case
%   Carrier2002_4 -> carrier2002Incommensurability % Anchor: incommensurability + empirical comparability
%   Malone1993_5  -> malone1993KuhnReconstructed  % Anchor: standard reading of incommensurability-without-relativism
%   Shimony1976_10 -> shimony1976KuhnTheses       % Anchor: classical philosophical engagement with Kuhn
%   Siegel1987_15 -> siegel1987KuhnInspired       % Anchor: survey of post-Kuhnian philosophy
%   Brewer2010_14 -> brewerLoschky2010TopDown     % Anchor: cognitive-psychology evidence for theory-laden observation
%   Lewowicz2011_2 -> lewowicz2011PurloinedReferent % Phlogiston referent
%   Boantza2008_4 -> boantza2008PhlogisticHeat    % Phlogiston-side defence
%   Halloun2016_8 -> halloun2016ModelLadenInquiry % Model-laden framing in pedagogy

\section{Theory-Laden Perception}
\label{sec:theoryladen}

\paragraph{Kuhn's three faces.}
The central phenomenological commitment of this paper is the
Kuhnian observation that the paradigm a scientist holds shapes what
she is even capable of perceiving \citep{kuhn1962Structure}. This is
not a methodological lapse to be corrected by more careful
experimental design but a structural feature of scientific cognition,
and it has three interconnected faces that we will treat as
expressions of a single underlying phenomenon. The first face is
what a scientist sees in incoming data: the same observation is read
differently from inside different paradigms, with the assignment of
\emph{anomaly} versus \emph{nuisance} itself being paradigm-dependent
\citep{Brewer2010_14, Shimony1976_10}. The second face is what
experiments she chooses to run: the paradigm shapes the space of
investigations she finds worth undertaking, which in turn shapes
what evidence the community as a whole produces
\citep{Halloun2016_8}. The third face is what theoretical
alternatives she even considers: rival frameworks are not so much
rejected as never seriously entertained, with the candidate set of
live revisions itself being structured by the prevailing framework
\citep{Carrier2002_4, Malone1993_5}.

\paragraph{Theory-ladenness as structural feature, not bias.}
The contemporary literature has substantially clarified what was, in
Kuhn's original framing, sometimes left ambiguous. \citet{Siegel1987_15}
and \citet{Malone1993_5} reconstruct incommensurability as a
condition on referential and inferential overlap between frameworks
rather than as a relativist denial of cross-paradigm rationality, and
\citet{Carrier2002_4} demonstrates on the phlogiston case that
\emph{empirical comparability} between rival paradigms is possible
without their being intertranslatable term-by-term. Cognitive
psychology has independently corroborated the perceptual layer of the
claim: \citet{Brewer2010_14} review top-down and bottom-up
contributions to scientific observation and find that the prior model
substantively conditions what counts as a datum, in line with the
phenomenology Kuhn was describing. The resulting picture is
deflationary in one direction --- theory-ladenness is not the
collapse of objectivity --- and inflationary in another: the
influence of the prevailing framework reaches further into the
cognitive pipeline than an additive bias term would suggest, touching
attention, salience, the felt space of alternatives, and the
allocation of experimental effort.

\paragraph{The Chemical Revolution as a worked case.}
The Chemical Revolution is the canonical historical demonstration of
all three faces operating jointly. \citet{Chang2012_9} argues that
phlogiston theory was not a failed predecessor refuted by Lavoisier's
data but a research programme with its own coherent set of
observables, anomalies, and live revisions; its replacement was less
a victory of evidence over confusion than a reallocation of what
counted as a phenomenon to be explained.
\citet{Lewowicz2011_2} and \citet{Boantza2008_4} reinforce this
reading from inside the chemistry: the phlogiston framework had
resources --- including a treatment of heat
\citep{Boantza2008_4} --- whose referents were not so much falsified
as relocated under the oxygen framework, exactly the kind of move
that the third Kuhnian face predicts. This historical record matters
for our model because it shows that the three faces are not
independent variables to be added: shifts in what is observed, what
is investigated, and what is entertained co-vary, and they co-vary
on the timescale of community-wide reorganisation rather than
individual reasoning.

\paragraph{Methodological commitment.}
We take this as the empirical commitment that licenses the model's
complexity. An accumulated-checklist alternative --- enumerate
biases in perception, biases in experiment choice, and biases in
hypothesis generation, then sum them --- would be both more
defensive and less informative, because it would treat the three
faces as separable failure modes rather than as expressions of a
single structural fact about how a community-held framework
conditions cognition. The stronger rhetorical position, and the one
this paper takes, is that a single underlying mechanism is
responsible for all three faces, and that the model must therefore
couple perception, action, and hypothesis-space in a way that an
additive bias account cannot. The remainder of the paper develops
that mechanism. In our setting, the
shared framework is encoded in the joint structure of agents'
generative models and the channel of inter-agent precision that
binds them; the three faces will reappear, respectively, as the
likelihood-weighting in inference, as the experiment-selection
policy, and as the support of the prior over theories. The case
for treating them together, rather than enumerating them
separately, rests on the phenomenological commitment defended
here.
